\documentclass[11pt]{article}
\usepackage{tikz}
\usetikzlibrary{arrows.meta,calc,positioning,patterns}

% Global preamble definitions
\definecolor{accentA}{RGB}{40,110,210}
\definecolor{accentB}{RGB}{220,90,40}
\tikzset{
  baseNode/.style={draw, rounded corners=1pt, inner sep=3pt, font=\small},
  flowArrow/.style={->, line width=0.4pt},
  helperGrid/.style={gray!30, line width=0.2pt}
}
\newcommand{\NodeStyle}[1]{baseNode, fill=#1!12}
\newcommand{\Tick}[1]{\node[font=\scriptsize] at #1 {};}
\newcommand{\Vec}[1]{\mathbf{#1}}

\begin{document}
\section*{Example Paper}
This fixture intentionally mixes prose, equations, and multiple TikZ pictures.

\begin{tikzpicture}[x=1cm,y=1cm]
  \draw[helperGrid] (-0.2,-0.2) grid (4.2,2.2);
  \node[\NodeStyle{accentA}] (A) at (0.8,1.5) {Input};
  \node[\NodeStyle{accentB}, right=1.4cm of A] (B) {Transform};
  \node[\NodeStyle{accentA}, right=1.4cm of B] (C) {Output};
  \draw[flowArrow] (A) -- (B);
  \draw[flowArrow] (B) -- (C);
\end{tikzpicture}

Some discussion between figures. We reference $\Vec{x}_{t+1}=A\Vec{x}_t+b$.

% Between-figure redefinitions
\renewcommand{\Vec}[1]{\boldsymbol{#1}}
\tikzset{
  flowArrow/.style={->, line width=0.55pt},
  emphasisBox/.style={draw=accentB, thick, rounded corners=2pt, fill=accentB!8}
}
\def\PlotScale{1.2}

\begin{tikzpicture}[scale=\PlotScale]
  \draw[->] (-0.2,0) -- (3.2,0) node[right] {$x$};
  \draw[->] (0,-0.2) -- (0,2.2) node[above] {$y$};
  \draw[accentA, thick, domain=0:3, samples=60] plot (\x,{0.2 + 0.6*\x - 0.1*\x*\x});
  \node[emphasisBox] at (2.2,1.7) {Updated style};
\end{tikzpicture}

More text. We now redefine local helper behavior before the next drawing.

% Additional between-figure macros
\newcommand{\StateNode}[2]{\node[\NodeStyle{accentA}] (#1) at #2 {#1};}
\renewcommand{\NodeStyle}[1]{baseNode, fill=#1!20, minimum width=8mm}
\tikzstyle{edgeNote}=[font=\scriptsize, fill=white, inner sep=1pt]

\begin{tikzpicture}
  \StateNode{S0}{(0,0)}
  \StateNode{S1}{(1.8,0.8)}
  \StateNode{S2}{(3.6,0)}
  \draw[flowArrow] (S0) -- node[edgeNote, above] {$a$} (S1);
  \draw[flowArrow] (S1) -- node[edgeNote, above] {$b$} (S2);
  \draw[flowArrow] (S2) to[bend left=22] node[edgeNote, below] {$c$} (S0);
\end{tikzpicture}

Final paragraph before the fourth figure.

% Last between-figure change
\tikzset{
  region/.style={pattern=north east lines, pattern color=accentA!45, draw=accentA},
  marker/.style={circle, fill=accentB, inner sep=1.1pt}
}
\let\OldTick\Tick
\renewcommand{\Tick}[1]{\node[marker] at #1 {};}

\begin{tikzpicture}[x=0.9cm,y=0.9cm]
  \draw[->] (0,0) -- (5,0);
  \draw[->] (0,0) -- (0,3.4);
  \draw[region] (0.8,0.6) rectangle (4.2,2.7);
  \Tick{(1.2,1.1)}
  \Tick{(2.0,2.1)}
  \Tick{(3.4,1.4)}
  \node[baseNode, fill=white] at (2.5,3.0) {Region of interest};
\end{tikzpicture}

\end{document}
